% 子文件：F_变换_初中
% 本文件包含 14 个 section


%% ==============================
%% 第7题：方格纸中的图形旋转、面积计算与格点确定
%% ==============================
\newpage

%% ==============================
%% 第7题：方格纸中的图形旋转、面积计算与格点确定
%% ==============================
\newpage

\section{解答：方格纸中的图形旋转、面积计算与确定格点}

\vspace{0.6cm}

\textbf{（1）以点 D 原点作中心对称画出 $180^\circ$ 旋转后的 $\Delta A_1B_1C_1$，以及（3）在直线 $BC$ 上确定格点 $E$ ，具体绘图解答如下：}

\vspace{0.4cm}

\begin{center}
\begin{tikzpicture}[scale=0.6, >=Stealth, line join=round, line cap=round]
  % 1. 绘制背景网格系 (12列 x 10行)
  % 我们以旋转中心 D(0,0) 为原点建立笛卡尔直角坐标系
  \draw[dashed, gray!80, step=1] (-6, -6) grid (6, 5); % 向下拓展到 -6 显示全 C1 和 B1 的对称图形需求
  
  % 2. 定义原三角形顶点及 D 点的位置坐标
  \coordinate (A) at (2, 4);
  \coordinate (B) at (-3, 4);
  \coordinate (C) at (5, 0);
  \coordinate (D) at (0, 0);
  
  % 标记原始点 D
  \fill (D) circle (2.5pt) node[below left] {$D$};
  
  % 绘制已给的黑色实线原三角形 ABC
  \draw[line width=1.0pt, black] (A) -- (B) -- (C) -- cycle;
  \node[above] at (A) {$A$};
  \node[above left] at (B) {$B$};
  \node[right] at (C) {$C$};
  
  % ----------------------------------------------------
  % 解答问题 (1)：画出旋转 180° 后的三角形 A1B1C1
  % ----------------------------------------------------
  % 中心对称旋转 180°，坐标关于原点中心取反 (-x, -y)
  \coordinate (A1) at (-2, -4);
  \coordinate (B1) at (3, -4);
  \coordinate (C1) at (-5, 0);
  
  % 绘制新的像三角形 A1B1C1, 使用蓝色标示特征作图连线
  \draw[line width=1.2pt, blue!80!black] (A1) -- (B1) -- (C1) -- cycle;
  \fill[blue!80!black] (A1) circle (2.5pt) node[below] {$A_1$};
  \fill[blue!80!black] (B1) circle (2.5pt) node[below right] {$B_1$};
  \fill[blue!80!black] (C1) circle (2.5pt) node[above left] {$C_1$};
  
  % 画出对应的原相-映像对称验证穿中心虚线
  \draw[dashed, gray, line width=0.4pt] (A) -- (A1);
  \draw[dashed, gray, line width=0.4pt] (B) -- (B1);
  \draw[dashed, gray, line width=0.4pt] (C) -- (C1);
  
  % ----------------------------------------------------
  % 解答问题 (3)：在直线 BC 上确定点 E，使得 BC = 2BE
  % ----------------------------------------------------
  % 根据已知 B(-3,4)，C(5,0)，使 BC=2BE 的正解即为其中点坐标 E(1,2)
  \coordinate (E) at (1, 2);
  \fill[red!80!black] (E) circle (3pt) node[above right] {$E$};
  \draw[dashed, red!80!black, line width=1.0pt] (B) -- (E);

  % ----------------------------------------------------
  % 附加面积求解平行四边形：辅助色图层
  % ----------------------------------------------------
  \fill[yellow!30, opacity=0.4] (B) -- (C1) -- (B1) -- (C) -- cycle;
  \draw[thick, dotted, orange] (B) -- (C1) -- (B1) -- (C) -- cycle;
  
\end{tikzpicture}
\end{center}

\vspace{0.6cm}

\textbf{（2）以 $B, C_1, B_1, C$ 为顶点的四边形的面积为} \underline{\quad \textbf{40} \quad} \textbf{。}

\vspace{0.6cm}

{\small\color{gray}
\textbf{解题说明：}\\
1. \textbf{第（1）小题：}\\
建立以中心点 $D(0,0)$ 为原点的直角坐标系。结合题意参数获取顶点坐标为 $A(2,4), B(-3,4), C(5,0)$。沿原点中心对称即为横纵坐标分别取反：$x \rightarrow -x, y \rightarrow -y$。由此可得出对称图像的顶点：$A_1(-2,-4), B_1(3,-4), C_1(-5,0)$。顺次连接这三点即得到所需图示 $\Delta A_1B_1C_1$。\\
2. \textbf{第（2）小题：}\\
由图形旋转性质可知由于对称性，该四边形的对角线平分交汇于点 $D$，可知四边形 $BC_1B_1C$ 必定属于平行四边形。其四个顶点为 $B(-3,4)$、$C_1(-5,0)$、$B_1(3,-4)$、$C(5,0)$。\\
利用经典的外接矩形割补法求面积：构造包裹该四边形的最小正外接矩形，可知其横轴长为 $5 - (-5) = 10$，纵轴高为 $4 - (-4) = 8$，全矩形总框面积为 $10 \times 8 = 80$。将四个角的“补全直角三角形”面积逐次减去，其大小依对称性分别相加得：$S_\text{外减} = 2 \times (\frac{1}{2} \times 2 \times 4 + \frac{1}{2} \times 8 \times 4) = 2 \times (4 + 16) = 40$。最终纯四边形净面积为 $80 - 40 = 40$。\\
3. \textbf{第（3）小题：}\\
由 $B(-3,4)$ 和 $C(5,0)$ 落定直线。在直线 $BC$ 上有格点 $E$，满足距离比例约束方程 $BC=2BE$。据此其物理含义为：点 $B$ 到 $E$ 的线段长度等于总线段 $BC$ 的二分之一。我们利用中点距离坐标公式：$E$ 的横坐标为 $\frac{-3+5}{2} = 1$，纵坐标为 $\frac{4+0}{2}=2$。得到解 $E(1, 2)$。其横纵属性均为整数，即符合“格点”要求，由此标明位置结束制图。
}

\vspace{0.6cm}

{\small\color{blue!70!black}
\textbf{绘图原理：}\\
\textbf{1. 网格化建立：}以点 $D$ 为直角坐标原点 $(0,0)$，根据设定参数直接通过 \texttt{step=1} 以及 \texttt{(-6, -6) grid (6, 5)} 搭建宽幅适配全对称旋转画面的底层虚线格坐标背景板。\\
\textbf{2. 代数原相作图：}依据人工输入的绝对单位节点坐标定义所有核心矢量定点图元，使用 TikZ \texttt{-{}-{}} 指令相连接完成原形 $\Delta ABC$，并计算赋入取反坐标画出高对比度蓝色填色的中心对称影像 $\Delta A_1B_1C_1$ 图层，保障线框逻辑稳定无视觉误差。\\
\textbf{3. 数学解析图视化加盖：}对于格点要求寻找，调用绝对极坐标及公式抓取点位算出唯一网格交点标上红底点缀。并辅设平行四边形的外轮廓面透白着色渲染 \texttt{opacity=0.4}，呈现了题目隐式存在的代数方程几何切割结果。
}

%% ==============================
%% 第8题：在方格纸中作图形的中心对称与复合变换
%% ==============================
\newpage



%% ==============================
%% 第8题：在方格纸中作图形的中心对称与复合变换
%% ==============================
\newpage

\section{解答：作图形的中心对称与复合变换图解}

\vspace{0.6cm}

\textbf{（1）画出 $\Delta A'B'C'$，及（2）进行图形变换的画图证明如下：}

\vspace{0.4cm}

\begin{center}
\begin{tikzpicture}[scale=0.7, >=Stealth, line join=round, line cap=round]
  % 1. 绘制背景网格系
  % 原题虽为 10x10，但在实际中心对称及后续平移中，图形在 y 轴向下穿透了底层边界。
  % 因此我们建立横竖更为宽裕的 10x12 拓展网格系，完全容纳所有解题辅助线，x: -5到5，y: -6到6
  \draw[dashed, gray!80, step=1] (-5, -6) grid (5, 6);
  
  % 2. 标记旋转中心 O
  \coordinate (O) at (0, 0);
  \fill (O) circle (2.5pt) node[above right] {$O$};
  
  % 3. 定义原始三角形 ABC
  \coordinate (A) at (-1, 2);
  \coordinate (B) at (-3, 5);
  \coordinate (C) at (-3, 2);
  
  % 绘制原始三角形
  \draw[line width=1.0pt, black] (A) -- (B) -- (C) -- cycle;
  \node[right] at (A) {$A$};
  \node[above left] at (B) {$B$};
  \node[below left] at (C) {$C$};
  
  % ----------------------------------------------------
  % 第一题：画出绕点 O 的中心对称图形 A'B'C'
  % ----------------------------------------------------
  % 中心对称相当于原点直接取反 (-x, -y)
  \coordinate (A1) at (1, -2);
  \coordinate (B1) at (3, -5);
  \coordinate (C1) at (3, -2);
  
  % 绘制对称像 A'B'C' (加粗蓝色突出答案)
  \draw[line width=1.5pt, blue!80!black] (A1) -- (B1) -- (C1) -- cycle;
  \fill[blue!80!black] (A1) circle (2pt) node[below left] {$A'$}; 
  \fill[blue!80!black] (B1) circle (2pt) node[right] {$B'$};
  \fill[blue!80!black] (C1) circle (2pt) node[above right] {$C'$};
  
  % 穿过原点 O 的中心对称连线点痕迹 (虚线)
  \draw[dashed, gray, line width=0.6pt] (A) -- (A1);
  \draw[dashed, gray, line width=0.6pt] (B) -- (B1);
  \draw[dashed, gray, line width=0.6pt] (C) -- (C1);

  % ----------------------------------------------------
  % 第二题：复合变换路径图形证明 (平移 + 旋转组合)
  % ----------------------------------------------------
  % 步骤I：将 ABC 整体竖直向下平移 7 个单位，得到基准点过渡图形 A''B''C''
  \coordinate (A2) at (-1, -5);
  \coordinate (B2) at (-3, -2);
  \coordinate (C2) at (-3, -5);
  
  % 绘制第一步中间态过渡图形 (黑色加粗虚线)
  \draw[dashed, black, line width=1.2pt] (A2) -- (B2) -- (C2) -- cycle;
  
  % 绘制平移动画箭头符号 (从各顶点向下垂直发射)
  \draw[->, dashed, red!70!black, line width=0.8pt] (B) -- (B2);
  \draw[->, dashed, red!70!black, line width=0.8pt] (C) -- (C2);
  \draw[->, dashed, red!70!black, line width=0.8pt] (A) -- (A2);
  

  % 步骤II：找出新的绝对对称中心点 O''。将过渡图形 A''B''C'' 绕新点作 180度 中心旋转，扣合 A'B'C'
  % 中点为 A2 与 A1 的中点：x = 0, y = -3.5
  \coordinate (O2) at (0, -3.5);
  \fill[red!80!black] (O2) circle (2.5pt) node[above] {$O''$};

  
  % 绘制二次旋转所穿过交汇于 O'' 的辅助线
  \draw[dashed, red!60!black, line width=0.6pt] (A2) -- (A1);
  \draw[dashed, red!60!black, line width=0.6pt] (B2) -- (B1);
  \draw[dashed, red!60!black, line width=0.6pt] (C2) -- (C1);
  
\end{tikzpicture}
\end{center}

\vspace{0.6cm}

\textbf{（2）解答：\underline{\quad 能 \quad}。结合作图说明如图及以下步骤解析。}

\vspace{0.6cm}

{\small\color{gray}
\textbf{解题说明：}\\
1. \textbf{第（1）问：中心对称作图}\\
以 $O$ 为原点构建格点坐标系，读出原始三个坐标 $A(-1,2), B(-3,5), C(-3,2)$。中心对称点即是在原坐标前加负号（等价于绕点旋转 $180^\circ$）。经计算得对应点 $A'(1,-2), B'(3,-5), C'(3,-2)$。描这些点并用实线顺次连接即可得到由深色标出的最终 $\Delta A'B'C'$。\\
\textbf{【提醒】}：严格依据原始的 $10\times10$ 格子，对称生成的底端点 $B'$ 实际上会越出底格下边界 $1$ 个单位，因此这里的作图辅助格已经动态向下扩大，使其能完整容纳图形展示！\\
2. \textbf{第（2）问：图形两次复合变换路径设计}\\
根据题意探究是否能用两次**不同的**图形变换（例如：平移、旋转、或者轴对称）的接续组合来实现目标效果。我们仔细观察图形相对位置，完全可以规划出完美干净的“竖直向下平移 $\rightarrow$ 中心交叉旋转”两步复合逻辑：\\
\textbf{步骤一：}将原始 $\Delta ABC$ 竖直向下平移 $7$ 个格点单位，让图形来到第三象限下方的位置。该过渡位置记为辅助虚线框 $\Delta A''B''C''$（其定点落子在 $A''(-1,-5), B''(-3,-2), C''(-3,-5)$）。\\
\textbf{步骤二：}接着将这个平移落位的过渡图形 $\Delta A''B''C''$ 进行旋转；我们取它与最终目标图形 $\Delta A'B'C'$ 处于镜像位置的两点连线（如 $C''$ 连到 $C'$）并取中点。经计算该中点坐标为新对称中心 $O''(0, -3.5)$。只需将过渡图形 $\Delta A''B''C''$ 借位绕点 $O''$ 中心对称旋转 $180^\circ$，这就通过严丝合缝的两次截然不同性质的变换（一为平移，一为旋转）获得了答案所指的 $\Delta A'B'C'$。
}

\vspace{0.6cm}

{\small\color{blue!70!black}
\textbf{绘图基本原理与步骤：}\\
\textbf{1. 突破画布自适应系建立：}完全抛弃由于母题 $10\times10$ 太过紧凑导致答卷超限的图形撕裂状况，程序建立以 $O(0,0)$ 原点的延伸网格。调用 \texttt{(-5, -6) grid (5, 6)} 全尺寸构建出更为广阔合理的 $10\times12$ 绝对格子阵列作为推演依托底层。\\
\textbf{2. 矩阵正切映射描图：}针对主像目标进行负号全导数转化计算高亮色块蓝线；并且对于题（2）隐含的二次变换：我们通过定义了一整套虚线外扩渲染样式，用强方向感 \texttt{->} 红色虚线将纵向无变形下的平移矢量轨迹“一次平移”呈现，紧缩算出虚拟交叉中点 $O''(0, -3.5)$ ，画出穿芯虚线辐射组，将“二次旋转”的抽象行为转化成了极度真实的纯物理结构图说明，全角度展示几何逻辑闭环。
}



%% ==============================
%% 第22题：一种花布售价与长度成正比例图像作图与填表
%% ==============================


%% ==============================
%% 第32题（补充题）：平移组合图案设计
%% ==============================
\newpage

\section{解答：平移组合图案设计}

\vspace{0.6cm}

\textbf{\LARGE 17.} \textbf{设计一个图案，使其可以由一个基本图形经过多次平移后组合得到，并给它赋予一定的含义。}

\vspace{0.4cm}
\textbf{解：}

\textbf{（1）设计思路：}

选取\textbf{"拼图块"}（带凸凹卡口的异形方块）作为基本图形。该形状的右侧有一个半圆形凸起，左侧有一个对应的半圆形凹槽：
\begin{itemize}
    \item \textbf{水平平移：}将基本图形向右平移 $1.8$ 厘米，后一块的左侧凹槽恰好吞入前一块的右侧凸起，形成无缝咬合。重复 $3$ 次，组成一行 $4$ 块互锁链条。
    \item \textbf{垂直平移：}将整行向下平移 $1.4$ 厘米，重复 $2$ 次，形成 $3$ 行 $\times$ $4$ 列共 $12$ 块全方位镶嵌阵列。
\end{itemize}

\textbf{图案含义：} \lq \lq 团结协作\rq \rq ——每一块拼图都不可或缺，只有通过紧密配合（平移拼接），才能组成完整的画面，象征团队合作的力量。

\vspace{0.4cm}
\textbf{（2）基本图形与完整图案：}

\vspace{0.3cm}
\begin{center}
\begin{tikzpicture}[>=Stealth, line join=round, line cap=round]

  % ---- 定义拼图块路径宏 ----
  % 参数 #1=x偏移, #2=y偏移, #3=填充色
  % 块体尺寸：宽约1.8cm，高约1.4cm
  % 右侧凸起半圆（半径0.18），左侧凹入半圆（半径0.18）
  \newcommand{\puzzle}[3]{
    \begin{scope}[xshift=#1 cm, yshift=#2 cm]
      \fill[#3, draw=black!70, line width=0.5pt]
        (0, 0) -- (0.7, 0)
        % 底部凸起
        arc[start angle=180, end angle=0, radius=0.18]
        -- (1.8, 0)
        -- (1.8, 0.5)
        % 右侧凸起
        arc[start angle=-90, end angle=90, radius=0.18]
        -- (1.8, 1.4)
        -- (1.06, 1.4)
        % 顶部凹入
        arc[start angle=0, end angle=180, radius=0.18]
        -- (0, 1.4)
        -- (0, 0.86)
        % 左侧凹入
        arc[start angle=90, end angle=-90, radius=0.18]
        -- cycle;
    \end{scope}
  }

  % ---- 左侧：展示基本图形 ----
  \node[font=\large\bfseries, above] at (0.9, 2.0) {基本图形};
  \puzzle{0}{0}{orange!40}

  % 标注凸凹特征
  \draw[->, line width=0.5pt, blue!80!black] (2.3, 0.7) -- (2.0, 0.7);
  \node[right, font=\footnotesize, text=blue!80!black] at (2.3, 0.7) {凸起};
  \draw[->, line width=0.5pt, blue!80!black] (-0.6, 0.7) -- (-0.15, 0.7);
  \node[left, font=\footnotesize, text=blue!80!black] at (-0.6, 0.7) {凹槽};

  % ---- 中间：平移箭头指引 ----
  \draw[->, line width=1.2pt, blue!80!black] (3.4, 0.7) -- (4.6, 0.7);
  \node[above, font=\small, text=blue!80!black] at (4, 0.8) {平移};

  % ---- 右侧：完整图案（4列 × 3行拼图阵列）----
  \node[font=\large\bfseries, above] at (8.1, 5.0) {完整图案：\lq \lq 团结协作\rq \rq };

  % 定义四种颜色交替排布
  % 第 1 行 (y=2.8)
  \puzzle{4.5}{2.8}{red!30}
  \puzzle{6.3}{2.8}{cyan!30}
  \puzzle{8.1}{2.8}{yellow!40}
  \puzzle{9.9}{2.8}{green!30}

  % 第 2 行 (y=1.4)
  \puzzle{4.5}{1.4}{cyan!30}
  \puzzle{6.3}{1.4}{yellow!40}
  \puzzle{8.1}{1.4}{green!30}
  \puzzle{9.9}{1.4}{red!30}

  % 第 3 行 (y=0)
  \puzzle{4.5}{0}{yellow!40}
  \puzzle{6.3}{0}{green!30}
  \puzzle{8.1}{0}{red!30}
  \puzzle{9.9}{0}{cyan!30}

  % ---- 标注平移方向与距离 ----
  % 水平平移标注
  \draw[<->, line width=0.6pt, blue!80!black] (5.4, 4.6) -- (7.2, 4.6);
  \node[above, font=\footnotesize, text=blue!80!black] at (6.3, 4.6) {向右平移};

  % 垂直平移标注
  \draw[<->, line width=0.6pt, blue!80!black] (12.2, 2.8) -- (12.2, 1.4);
  \node[right, font=\footnotesize, text=blue!80!black] at (12.2, 2.1) {向下平移};

\end{tikzpicture}
\end{center}

\vspace{0.4cm}
{\small\color{blue!70!black}
\textbf{绘图原理与步骤：}\\
\textbf{1. 基本图形构造：} 拼图块主体为 $1.8 \times 1.4$ cm 矩形，右侧和底部各带一个半径 $0.18$ cm 的半圆凸起（\texttt{arc} 外弧），左侧和顶部各带一个对应的半圆凹槽（\texttt{arc} 内弧）。凸凹尺寸严格匹配，确保平移后无缝咬合。\\
\textbf{2. 平移实现方式：} 定义 \texttt{\textbackslash{}puzzle\{x\}\{y\}\{color\}} 宏，通过 \texttt{xshift/yshift} 控制平移。水平间距 $1.8$ cm（等于块宽），垂直间距 $1.4$ cm（等于块高），共 $4 \times 3 = 12$ 块。每块凸起恰好嵌入相邻块的凹槽。\\
\textbf{3. 颜色编码：} 使用红、青、黄、绿四色循环填充，使相邻块易于区分，清晰展示每块的独立性与整体的组合性。\\
\textbf{4. 标注系统：} 左侧展示单个拼图块并标注凸凹特征，右侧完整阵列附带平移方向和距离标注。
}
%% ==============================
%% 第33题（补充题）：多边形轴对称图形判断
%% ==============================
\newpage



%% ==============================
%% 第58题（补充题）：网格中的图形平移
%% ==============================
\newpage

\section{解答：网格图中的图形平移}

\vspace{0.6cm}

\textbf{\LARGE 49.} \textbf{（例 1）如图 9-2，平移 $\triangle ABC$，使点 $A$ 移动到点 $D$ 的位置。}

\textbf{（1）画出平移后所得的 $\triangle DEF$，并写出对应点、对应边、对应角。}

\textbf{（2）结合平移的定义，写出 $\triangle DEF$ 是由 $\triangle ABC$ 经过怎样的平移得到的。}

\vspace{0.4cm}
\textbf{解：}

\textbf{（1）画图如下：}

\begin{center}
\begin{tikzpicture}[>=Stealth, scale=0.8]
  % 绘制网格背景 (13x7)
  \draw[dashed, gray] (0,0) grid (15,9);
  
  % 定义坐标
  \coordinate (A) at (5,7);
  \coordinate (B) at (3,3);
  \coordinate (C) at (7,4);
  \coordinate (D) at (13,6);
  \coordinate (E) at (11,2);
  \coordinate (F) at (15,3);
  
  % 绘制原三角形 ABC (黑色)
  \draw[thick] (A) -- (B) -- (C) -- cycle;
  \filldraw[black] (A) circle (1.5pt);
  \filldraw[black] (B) circle (1.5pt);
  \filldraw[black] (C) circle (1.5pt);
  \node[above] at (A) {$A$};
  \node[below left] at (B) {$B$};
  \node[below right] at (C) {$C$};
  
  % 绘制平移后的三角形 DEF (红色显示区分)
  \draw[thick, red] (D) -- (E) -- (F) -- cycle;
  \filldraw[red] (D) circle (1.5pt);
  \filldraw[red] (E) circle (1.5pt);
  \filldraw[red] (F) circle (1.5pt);
  \node[above right, text=red] at (D) {$D$};
  \node[below left, text=red] at (E) {$E$};
  \node[below right, text=red] at (F) {$F$};

  % 补充平移的虚线箭头
  \draw[->, dashed, blue!70] (A) -- (D);
  \draw[->, dashed, blue!70] (B) -- (E);
  \draw[->, dashed, blue!70] (C) -- (F);

\end{tikzpicture}
\end{center}

\vspace{0.2cm}
\textbf{对应关系：}

对应点：点 $A$ 和点 $D$，点 $B$ 和点 $E$，点 $C$ 和点 $F$。

对应边：线段 $AB$ 和 $DE$，线段 $BC$ 和 $EF$，线段 $AC$ 和 $DF$。

对应角：$\angle BAC$ 和 $\angle EDF$，$\angle ABC$ 和 $\angle DEF$，$\angle ACB$ 和 $\angle DFE$。

\vspace{0.4cm}
\textbf{（2）平移过程说明：}

观察网格可知，把点 $A$ 移动到点 $D$ 的位置，相当于将其\textbf{向右平移了 $\mathbf{6}$ 个单位长度，再向下平移了 $\mathbf{1}$ 个单位长度}。

因此，$\triangle DEF$ 是由 $\triangle ABC$ 经过\textbf{向右平移 $\mathbf{6}$ 个单位长度，再向下平移 $\mathbf{1}$ 个单位长度}得到的。（或者描述为：沿射线 $AD$ 的方向，平移线段 $AD$ 的长度。）

\vspace{0.4cm}
{\small\color{blue!70!black}
\textbf{绘图原理与步骤：}\\
\textbf{1. 网格与坐标系：} 建立大小为 $13 \times 7$ 的 \texttt{dashed} 虚线网格系统，左下角定为 $(0,0)$。\\
\textbf{2. 顶点推算：} 观察原图设定点 $B(2,2)$，向右 $3$ 向上 $1$ 得 $C(5,3)$，从 $B$ 向右 $2$ 向上 $3$ 定出 $A(4,5)$。从 $A$ 向右 $6$ 向下 $1$ 得目标点 $D(10,4)$。\\
\textbf{3. 向量平移：} 求出平移向量 $\vec{v} = (6, -1)$，得到对应点 $E = B + \vec{v} = (8,1)$ 和 $F = C + \vec{v} = (11,2)$。\\
\textbf{4. 图形渲染：} 用黑色实线绘制 $\triangle ABC$；用红色描绘平移后的 $\triangle DEF$ 并附以蓝色虚线箭头指示对应点的平移轨迹方向，使作图清晰直观且便于区分原图与目标图。
}

%% ==============================
%% 第59题（补充题）：三角形平移的两种作图方法
%% ==============================
\newpage



%% ==============================
%% 第59题（补充题）：三角形平移的两种作图方法
%% ==============================
\newpage

\section{解答：三角形平移的两种作图方法}

\vspace{0.6cm}

\textbf{\LARGE 50.} \textbf{（例 1）已知 $\triangle ABC$ 的顶点 $A$ 平移到顶点 $D$，请用两种不同的方法，在图 9-5 中作出平移后的图形。}

\vspace{0.4cm}
\textbf{解：}

\textbf{（1）方法一：利用“经过平移，对应点所连的线段平行且相等”作图。}

\textbf{作法：}

① 分别过点 $B$、$C$ 作射线 $AD$ 的平行线；

② 在所作的平行线上沿 $A \to D$ 的方向分别截取 $BE = AD$，$CF = AD$，得到点 $B$ 的对应点 $E$ 和点 $C$ 的对应点 $F$；

③ 顺次连接 $D$、$E$、$F$。$\triangle DEF$ 即为所求作的图形。

\begin{center}
\begin{tikzpicture}[>=Stealth, scale=0.9]
  % 坐标设定
  \coordinate (A) at (1, 3);
  \coordinate (B) at (0, 0);
  \coordinate (C) at (2.5, 0.5);
  \coordinate (D) at (5, 2.5);
  
  \coordinate (E) at (4, -0.5);
  \coordinate (F) at (6.5, 0);

  % 原三角形
  \draw[thick] (A) -- (B) -- (C) -- cycle;
  \filldraw (A) circle (1.2pt) node[above] {$A$};
  \filldraw (B) circle (1.2pt) node[below left] {$B$};
  \filldraw (C) circle (1.2pt) node[below right] {$C$};
  
  % 点D
  \filldraw (D) circle (1.2pt) node[above] {$D$};
  
  % 连线对应点 (作法一的特征)
  \draw[dashed, blue!80, ->] (A) -- (D);
  \draw[dashed, blue!80, ->] (B) -- (E);
  \draw[dashed, blue!80, ->] (C) -- (F);

  % 目标三角形
  \draw[thick, red] (D) -- (E) -- (F) -- cycle;
  \filldraw[red] (E) circle (1.2pt) node[below left, text=red] {$E$};
  \filldraw[red] (F) circle (1.2pt) node[below right, text=red] {$F$};
\end{tikzpicture}
\end{center}

\vspace{0.4cm}
\textbf{（2）方法二：利用“经过平移，对应线段平行且相等”作图。}

\textbf{作法：}

① 过点 $D$ 作线段 $AB$ 的平行线，并在该线上沿 $A \to B$ 的方向截取 $DE = AB$，得到点 $B$ 的对应点 $E$；

② 过点 $D$ 作线段 $AC$ 的平行线，并在该线上沿 $A \to C$ 的方向截取 $DF = AC$，得到点 $C$ 的对应点 $F$；

③ 连接 $EF$。$\triangle DEF$ 即为所求作的图形。

\begin{center}
\begin{tikzpicture}[>=Stealth, scale=0.9]
  % 坐标设定
  \coordinate (A) at (1, 3);
  \coordinate (B) at (0, 0);
  \coordinate (C) at (2.5, 0.5);
  \coordinate (D) at (5, 2.5);
  
  \coordinate (E) at (4, -0.5);
  \coordinate (F) at (6.5, 0);

  % 原三角形
  \draw[thick] (A) -- (B) -- (C) -- cycle;
  \filldraw (A) circle (1.2pt) node[above] {$A$};
  \filldraw (B) circle (1.2pt) node[below left] {$B$};
  \filldraw (C) circle (1.2pt) node[below right] {$C$};
  
  % 点D与平移向量示意
  \filldraw (D) circle (1.2pt) node[above] {$D$};
  \draw[dashed, gray] (A) -- (D);

  % 作法二的特征：作平行边 (只呈现最终红边和少许延长引线)
  \draw[dashed, orange] (D) -- (E); 
  \draw[dashed, orange] (D) -- (F);

  % 目标三角形
  \draw[thick, red] (D) -- (E) -- (F) -- cycle;
  \filldraw[red] (E) circle (1.2pt) node[below left, text=red] {$E$};
  \filldraw[red] (F) circle (1.2pt) node[below right, text=red] {$F$};
\end{tikzpicture}
\end{center}

\vspace{0.4cm}
{\small\color{blue!70!black}
\textbf{绘图原理与步骤：}\\
\textbf{1. 图形定位：} 设置原图基准比例点 $A(1,3), B(0,0), C(2.5,0.5)$ 以及目标点 $D(5,2.5)$。计算平移向量为 $\vec{v} = D-A = (4,-0.5)$，由此求出 $E = B + \vec{v} = (4,-0.5)$，及 $F = C + \vec{v} = (6.5,0)$。\\
\textbf{2. 方法一特征绘图：} 核心为“对应点连线长短一致、互相平行”。用\texttt{dashed, blue!80, ->}蓝色虚线箭头连接对应顶点 $A \to D$、$B \to E$、$C \to F$ 直观展示每个顶点的位移过程。\\
\textbf{3. 方法二特征绘图：} 核心为“对应边平行且相等”。图中省略了连接其余顶点的连线（仅保留了灰色的 $AD$ 参照），通过绘制橙色虚线底层与红色的平移线段重叠，体现由已知点 $D$ 引出的平行线逻辑。\\
\textbf{4. 图形呈现：} 统一原图图形为黑色实线，平移构成的目标图形 $\triangle DEF$ 加上红色实线与端点以高亮区分。两图上下呼应体现“异曲同工”。
}

%% ==============================
%% 第60题（补充题）：轴对称图形的画法
%% ==============================
\newpage



%% ==============================
%% 第60题（补充题）：轴对称图形的画法
%% ==============================
\newpage

\section{解答：轴对称图形的画法}

\vspace{0.6cm}

\textbf{\LARGE 51.} \textbf{（第 3 题）如图，画出 $\triangle ABC$ 关于直线 $l$ 对称的 $\triangle A'B'C'$。}

\vspace{0.4cm}
\textbf{解：}

\textbf{画图如下：}

\begin{center}
\begin{tikzpicture}[>=Stealth, scale=1.2]
  % 定义对称轴 l
  \draw[thick] (0, -1.5) -- (0, 4) node[left] {$l$};

  % 定义原始三角形的顶点坐标
  \coordinate (A) at (0, 2.5);
  \coordinate (B) at (-2.5, -0.5);
  \coordinate (C) at (-0.5, 0);

  % 根据对称轴 x=0 计算对称点坐标
  \coordinate (A') at (0, 2.5); % A 点在对称轴上
  \coordinate (B') at (2.5, -0.5);
  \coordinate (C') at (0.5, 0);

  % 绘制原始三角形 ABC (黑色)
  \draw[thick] (A) -- (B) -- (C) -- cycle;
  \node[right] at (0.1, 2.6) {$A(A')$}; % A和A'重合
  \node[below left] at (B) {$B$};
  \node[below left] at (C) {$C$};

  % 绘制对称三角形 A'B'C' (红色高亮)
  \draw[thick, red] (A') -- (B') -- (C') -- cycle;
  \node[below right, text=red] at (B') {$B'$};
  \node[below right, text=red] at (C') {$C'$};

  % 绘制辅助虚线（垂直于对称轴）
  \draw[dashed, blue!70] (B) -- (B');
  \draw[dashed, blue!70] (C) -- (C');
  
  % 标注重合点A与顶点
  \filldraw (A) circle (1.2pt);
  \filldraw (B) circle (1.2pt);
  \filldraw (C) circle (1.2pt);
  \filldraw[red] (B') circle (1.2pt);
  \filldraw[red] (C') circle (1.2pt);
  
  % 垂直标记
  \draw[blue!70] (0, -0.5) ++(0, 0.15) -- ++(-0.15, 0) -- ++(0, -0.15);
  \draw[blue!70] (0, 0) ++(0, 0.15) -- ++(-0.15, 0) -- ++(0, -0.15);

\end{tikzpicture}
\end{center}

\vspace{0.4cm}
\textbf{作图说明：}

① 因为点 $A$ 在对称轴 $l$ 上，所以点 $A$ 的对称点 $A'$ 就是点 $A$ 自身；

② 分别过点 $B$、$C$ 向直线 $l$ 作垂线，并延长至点 $B'$、$C'$，使得点 $B'$、$C'$ 到直线 $l$ 的距离分别等于点 $B$、$C$ 到直线 $l$ 的距离（即 $l$ 为线段 $BB'$、$CC'$ 的垂直平分线）；

③ 连接 $A'B'$、$B'C'$、$C'A'$，所得 $\triangle A'B'C'$ 即为所求作的图形。

\vspace{0.4cm}
{\small\color{blue!70!black}
\textbf{绘图原理与步骤：}\\
\textbf{1. 坐标系与对称轴：} 选取 $Y$ 轴（$x=0$）作为对称轴 $l$。\\
\textbf{2. 顶点设定：} 根据原图比例观察，设定 $A(0, 2.5)$ 在轴上，$B(-2.5, -0.5)$，$C(-0.5, 0)$。\\
\textbf{3. 计算对称点：} 根据基于 $y$ 轴的轴对称变换原理（横坐标取相反数），计算出 $B'(2.5, -0.5)$，$C'(0.5, 0)$，$A'$ 与 $A$ 重合。\\
\textbf{4. 图形渲染：} 原 $\triangle ABC$ 用黑色实线绘制，目标 $\triangle A'B'C'$ 用红色实线绘制高亮区分。利用蓝色虚线展示 $BB'$ 和 $CC'$ 的对称及垂直平分特征，并在与轴的交点处作直角符号增强逻辑严谨性。
}

%% ==============================
%% 第61题（补充题）：网格中四边形的轴对称变换
%% ==============================
\newpage



%% ==============================
%% 第61题（补充题）：网格中四边形的轴对称变换
%% ==============================
\newpage

\section{解答：网格图中的轴对称变换}

\vspace{0.6cm}

\textbf{\LARGE 52.} \textbf{（第 4 题）如图，在方格纸中（每个小方格是边长均等的正方形）画出四边形 $ABCD$ 关于直线 $l$ 对称的四边形 $A'B'C'D'$。}

\vspace{0.4cm}
\textbf{解：}

\textbf{画图如下：}

\begin{center}
\begin{tikzpicture}[>=Stealth, scale=0.8]
  % 绘制网格背景 (-5 到 5, 0 到 10)
  \draw[dashed, gray] (-5,0) grid (5,10);
  
  % 定义并绘制对称轴 l
  \draw[thick, black] (0, -0.5) -- (0, 10.5) node[left] {$l$};

  % 定义原始四边形的顶点坐标
  \coordinate (A) at (-3, 9);
  \coordinate (B) at (-4, 6);
  \coordinate (C) at (-1, 6);
  \coordinate (D) at (-1, 8);

  % 根据对称轴 x=0 计算对称点坐标
  \coordinate (A') at (3, 9);
  \coordinate (B') at (4, 6);
  \coordinate (C') at (1, 6);
  \coordinate (D') at (1, 8);

  % 绘制辅助虚线（垂直于对称轴）
  \draw[dashed, blue!70] (A) -- (A');
  \draw[dashed, blue!70] (B) -- (B');
  \draw[dashed, blue!70] (C) -- (C');
  \draw[dashed, blue!70] (D) -- (D');

  % 绘制原始四边形 ABCD (黑色)
  \draw[thick] (A) -- (B) -- (C) -- (D) -- cycle;
  \filldraw (A) circle (1.5pt);
  \filldraw (B) circle (1.5pt);
  \filldraw (C) circle (1.5pt);
  \filldraw (D) circle (1.5pt);
  \node[above left] at (A) {$A$};
  \node[below left] at (B) {$B$};
  \node[below right] at (C) {$C$};
  \node[above right] at (D) {$D$};

  % 绘制对称四边形 A'B'C'D' (红色高亮)
  \draw[thick, red] (A') -- (B') -- (C') -- (D') -- cycle;
  \filldraw[red] (A') circle (1.5pt);
  \filldraw[red] (B') circle (1.5pt);
  \filldraw[red] (C') circle (1.5pt);
  \filldraw[red] (D') circle (1.5pt);
  \node[above right, text=red] at (A') {$A'$};
  \node[below right, text=red] at (B') {$B'$};
  \node[below left, text=red] at (C') {$C'$};
  \node[above left, text=red] at (D') {$D'$};

\end{tikzpicture}
\end{center}

\vspace{0.4cm}
\textbf{作图说明：}

① 找出原四边形 $ABCD$ 的各顶点；

② 借助网格线，分别找到点 $A$、$B$、$C$、$D$ 关于直线 $l$ 对称的点 $A'$、$B'$、$C'$、$D'$。因为网格的横线均垂直于直线 $l$，所以对应点直接向右平移原来距离即可（即各对应点到直线 $l$ 的格数相等）；

③ 顺次连接 $A'B'$、$B'C'$、$C'D'$、$D'A'$，所得的四边形 $A'B'C'D'$ 即为所求作的图形。

\vspace{0.4cm}
{\small\color{blue!70!black}
\textbf{绘图原理与步骤：}\\
\textbf{1. 网格与坐标系：} 建立大小为 $11 \times 9$ 的 \texttt{dashed} 虚线网格系统，横坐标跨度自 $x=-6$ 到 $x=5$ 以吻合原图布局。对称轴 $l$ 设定为 $Y$ 轴（$x=0$），长度上下延伸半个刻度。\\
\textbf{2. 顶点设定：} 观察原图中网格距离推算，确定原特征坐标为 $A(-3,8)$、$B(-4,4)$、$C(-1,4)$、$D(-1,7)$。\\
\textbf{3. 对称计算：} 根据 $Y$ 轴的轴对称横纵变换原理，求出对应镜像坐标 $A'(3,8)$、$B'(4,4)$、$C'(1,4)$、$D'(1,7)$。\\
\textbf{4. 图形渲染：} 原四边形 $ABCD$ 以黑色实线绘制，目标四边形 $A'B'C'D'$ 以红色实线描绘作高亮区分，顶点连线用蓝色虚线表现出网格间作对称图形特征线时的跨度关系。
}

%% ==============================
%% 第62题（补充题）：市民热线统计表与扇形统计图
%% ==============================
\newpage



%% ==============================
%% 第63题（补充题）：网格中的三角形旋转
%% ==============================
\newpage

\section{解答：网格图中的图形旋转}

\vspace{0.6cm}

\textbf{\LARGE 54.} \textbf{（第 3 题）如图，在由边长为 $1$ 个单位长度的小正方形组成的方格纸中，$\triangle ABC$ 的顶点分别在格点上，以点 $C$ 为旋转中心，将 $\triangle ABC$ 按顺时针方向旋转 $90^\circ$，请在方格纸中画出旋转后的 $\triangle A'B'C$。}

\vspace{0.4cm}
\textbf{解：}

\textbf{画图如下：}

\begin{center}
\begin{tikzpicture}[>=Stealth, scale=0.6]
  % 绘制网格背景 (0 到 13, 0 到 11)
  \draw[dashed, gray] (0,0) grid (13,11);
  
  % 定义原始三角形的顶点坐标
  \coordinate (A) at (0, 11);
  \coordinate (B) at (3, 8);
  \coordinate (C) at (2, 5);

  % 以 C 为中心，顺时针旋转 90 度计算新坐标
  % C(2,5), A(0,11) -> 向量 CA = (-2, 6) -> 顺时针旋转 90 度变为 (6, 2) -> A' = (8, 7)
  % B(3,8) -> 向量 CB = (1, 3) -> 顺时针旋转 90 度变为 (3, -1) -> B' = (5, 4)
  \coordinate (A') at (8, 7);
  \coordinate (B') at (5, 4);

  % 绘制原始三角形 ABC (黑色)
  \draw[thick] (A) -- (B) -- (C) -- cycle;
  \filldraw (A) circle (2pt) node[above] {$A$};
  \filldraw (B) circle (2pt) node[above right] {$B$};
  \filldraw (C) circle (2pt) node[below left] {$C$};

  % 绘制旋转后的三角形 A'B'C (红色高亮)
  \draw[thick, red] (A') -- (B') -- (C) -- cycle;
  \filldraw[red] (A') circle (2pt) node[above right, text=red] {$A'$};
  \filldraw[red] (B') circle (2pt) node[below right, text=red] {$B'$};

  % 绘制辅助虚线（展现旋转轨迹）
  \draw[dashed, blue!70] (C) -- (A');
  
  % 标示旋转角 (CA 到 CA') 与 垂直符号
  % atan2(6, -2) 约 108.4 度，atan2(2, 6) 约 18.4 度
  \draw[blue!70, ->] (C) ++(108.4:1.2) arc (108.4:18.4:1.2);
  \node[blue!70, font=\small] at (4, 6.5) {$90^\circ$};
  
  % 直角（垂直）符号
  \draw[blue!70, thick] (C) ++(108.4:0.5) -- ++(18.4:0.5) -- ++(108.4:-0.5);

\end{tikzpicture}
\end{center}

\vspace{0.4cm}
\textbf{作图说明：}

① 根据题意，旋转中心为点 $C$，所以点 $C$ 旋转后的位置保持不变（即点 $C'$ 与点 $C$ 重合）；

② 将线段 $CA$ 绕点 $C$ 顺时针旋转 $90^\circ$ 得到线段 $CA'$。在网格中体现为：将原方向“向左 $2$ 格、向上 $6$ 格”转化为“向右 $6$ 格、向上 $2$ 格”，从而确定对应点 $A'$ 的位置；

③ 将线段 $CB$ 绕点 $C$ 顺时针旋转 $90^\circ$ 得到线段 $CB'$。在网格中体现为：将原方向“向右 $1$ 格、向上 $3$ 格”转化为“向右 $3$ 格、向下 $1$ 格”，从而确定对应点 $B'$ 的位置；

④ 顺次连接 $A'B'$、$B'C$，所得的 $\triangle A'B'C$ 即为所求作的旋转图形。

\vspace{0.4cm}
{\small\color{blue!70!black}
\textbf{绘图原理与步骤：}\\
\textbf{1. 网格系统：} 根据已知条件建立大小为 $13 \times 11$ 的虚线方格坐标系，设定缩放比例为 \texttt{scale=0.6} 以兼顾排版尺寸展示。\\
\textbf{2. 坐标推算：} 以左下角为原点 $(0,0)$，确定原三角形顶点为 $A(0,11), B(3,8), C(2,5)$。顺时针旋转 $90^\circ$ 的向量法则为 $(x, y) \to (y, -x)$。根据基于中心 $C$ 的相对向量坐标，推算出 $A'(8, 7)$ 及 $B'(5, 4)$。\\
\textbf{3. 图形展示：} 使用红色实线突出显示目标旋转图形 $\triangle A'B'C$，同时保留原图形的黑色实线对比。\\
\textbf{4. 过程标注：} 绘制蓝色虚线标出旋转前后的对应边 $CA'$，并使用带箭头的 \texttt{arc} 圆弧标记出 $90^\circ$ 的极坐标角变化轨迹动作，让画法痕迹显得更加丰满易懂。
}

%% ==============================
%% 第64题（补充题）：四边形绕点旋转
%% ==============================
\newpage



%% ==============================
%% 第64题（补充题）：四边形绕点旋转
%% ==============================
\newpage

\section{解答：四边形绕点旋转作图}

\vspace{0.6cm}

\textbf{\LARGE 55.} \textbf{（第 3 题）如图，画出四边形 $ABCD$ 绕点 $O$ 按顺时针方向旋转 $120^\circ$ 所得到的图形。}

\vspace{0.4cm}
\textbf{解：}

\textbf{画图如下：}

\begin{center}
\begin{tikzpicture}[>=Stealth, scale=0.8]
  % 定义中心 O
  \coordinate (O) at (0, 0);
  
  % 定义四边形 ABCD 坐标 (基于原图比例)
  \coordinate (A) at (-4, 1);
  \coordinate (B) at (-4.5, -1);
  \coordinate (C) at (-1.5, -1);
  \coordinate (D) at (-2, 1);
  
  % 原四边形
  \draw[thick] (A) -- (B) -- (C) -- (D) -- cycle;
  \filldraw (A) circle (1.5pt) node[left] {$A$};
  \filldraw (B) circle (1.5pt) node[left] {$B$};
  \filldraw (C) circle (1.5pt) node[below] {$C$};
  \filldraw (D) circle (1.5pt) node[above] {$D$};
  
  % 旋转中心 O
  \filldraw (O) circle (1.5pt) node[right] {$O$};
  
  % 计算并绘制目标四边形 (在局部空间发生坐标系旋转)
  \begin{scope}[rotate around={-120:(O)}]
    \coordinate (A') at (-4, 1);
    \coordinate (B') at (-4.5, -1);
    \coordinate (C') at (-1.5, -1);
    \coordinate (D') at (-2, 1);
    
    \draw[thick, red] (A') -- (B') -- (C') -- (D') -- cycle;
    \filldraw[red] (A') circle (1.5pt) node[above right, text=red] {$A'$};
    \filldraw[red] (B') circle (1.5pt) node[above, text=red] {$B'$};
    \filldraw[red] (C') circle (1.5pt) node[left, text=red] {$C'$};
    \filldraw[red] (D') circle (1.5pt) node[below right, text=red] {$D'$};
  \end{scope}
  
  % 画出辅助线，体现 120 度旋转
  \draw[dashed, blue!70] (D) -- (O) -- (D');
  
  % 画 -120 度的角度弧圈指示
  \draw[blue!70, ->] (O) ++(153.4:0.8) arc (153.4:33.4:0.8);
  \node[blue!70, font=\small] at (93.4:1.2) {$120^\circ$};
\end{tikzpicture}
\end{center}

\vspace{0.4cm}
\textbf{作图说明：}

① 连接 $OA$、$OB$、$OC$、$OD$；

② 分别做 $\angle AOA' = \angle BOB' = \angle COC' = \angle DOD' = 120^\circ$（按顺时针方向），且截取使 $OA'=OA$、$OB'=OB$、$OC'=OC$、$OD'=OD$，依次得到原顶点 $A$、$B$、$C$、$D$ 的对应目标点 $A'$、$B'$、$C'$、$D'$；

③ 顺次连接 $A'B'$、$B'C'$、$C'D'$、$D'A'$，所得的四边形 $A'B'C'D'$ 即为所求作的顺时针旋转图形。

\vspace{0.4cm}
{\small\color{blue!70!black}
\textbf{绘图原理与步骤：}\\
\textbf{1. 坐标系选定：} 设定旋转中心 $O$ 为原点 $(0,0)$。根据图像比例近似设定梯形边界基准：赋予四边形顶点坐标 $A(-4, 1), B(-4.5, -1), C(-1.5, -1), D(-2, 1)$。\\
\textbf{2. 旋转变换：} 运用 \texttt{TikZ} 的 \texttt{scope} 并叠加局部坐标变换指令 \texttt{rotate around=\{-120:(O)\}}，直接生成绝对顺时针 $120^\circ$ 的映射图，此举规避了手算各类浮点数三角函数的代码复杂度与精度丢失。\\
\textbf{3. 过程连线：} 利用浅蓝色虚线连接了一组距离中心 $O$ 较近的关键点映射 $D \to O \to D'$，利用 \texttt{arc} 圆弧表现出动态位移走向，避免全部虚线导致画面拥挤。
}

%% ==============================
%% 第65题（补充题）：平移等价于两次旋转
%% ==============================
\newpage



%% ==============================
%% 第65题（补充题）：平移等价于两次旋转
%% ==============================
\newpage

\section{解答：平移的两次中心旋转等效变换}

\vspace{0.6cm}

\textbf{\LARGE 56.} \textbf{（第 5 题）如图，已知 $\triangle A'B'C'$ 是由 $\triangle ABC$ 经过平移得到的。 $\triangle A'B'C'$ 是否还可以看作由 $\triangle ABC$ 经过两次旋转得到？若能，请画出示意图；若不能，请说明理由。}

\vspace{0.4cm}
\textbf{解：}

\textbf{能看作是由两次旋转得到的。}

\textbf{理由与示意图作法：} 

根据几何原理，任意一个**平移**变换都可以分解为**连续两次 $180^\circ$ 的旋转（即中心对称）**。
我们可以在空间中任意提取一个中转核心 $O_1$，将 $\triangle ABC$ 绕点 $O_1$ 先旋转 $180^\circ$ 得到过渡图形 $\triangle A_1B_1C_1$；然后再在过渡图形对应顶点（如 $A_1$）与目标图形对应顶点（如 $A'$）之间取中点作为第二个中心 $O_2$，将 $\triangle A_1B_1C_1$ 绕着点 $O_2$ 再度旋转 $180^\circ$，即可成功翻折得出平移体 $\triangle A'B'C'$。

\textbf{示意图如下：}

\begin{center}
\begin{tikzpicture}[>=Stealth, scale=0.8]
  % 定义原三角形坐标 (依比例 800宽, 630横 x 510高 缩放构建, 缩小200倍)
  \coordinate (A) at (0.15, 2.05);
  \coordinate (B) at (-3, -0.5);
  \coordinate (C) at (1, -0.5);

  % 设定平移向量 v = (5, 1)，得出目标三角形
  \coordinate (A') at (5.15, 3.05);
  \coordinate (B') at (2, 0.5);
  \coordinate (C') at (6, 0.5);

  % 任意选取第一次旋转中心 O_1
  \coordinate (O1) at (0.5, -1.5);
  
  % 计算过渡三角形 A1, B1, C1 (绕 O1 旋转 180 度，即基于 O1 的镜像点)
  \coordinate (A1) at (0.85, -5.05);
  \coordinate (B1) at (4, -2.5);
  \coordinate (C1) at (0, -2.5);
  
  % 计算第二次旋转中心 O_2，为 A1 和 A' 的中点
  \coordinate (O2) at (3, -1);

  % 过渡三角形 (灰色虚线)
  \draw[dashed, gray, thick] (A1) -- (B1) -- (C1) -- cycle;
  \filldraw[gray] (A1) circle (1.5pt) node[below] {$A_1$};
  \filldraw[gray] (B1) circle (1.5pt) node[right] {$B_1$};
  \filldraw[gray] (C1) circle (1.5pt) node[left] {$C_1$};

  % 原三角形
  \draw[thick] (A) -- (B) -- (C) -- cycle;
  \filldraw (A) circle (1.5pt) node[above] {$A$};
  \filldraw (B) circle (1.5pt) node[left] {$B$};
  \filldraw (C) circle (1.5pt) node[right] {$C$};

  % 目标三角形
  \draw[thick, red] (A') -- (B') -- (C') -- cycle;
  \filldraw[red] (A') circle (1.5pt) node[above, text=red] {$A'$};
  \filldraw[red] (B') circle (1.5pt) node[left, text=red] {$B'$};
  \filldraw[red] (C') circle (1.5pt) node[right, text=red] {$C'$};

  % 旋转中心点 (橙色)
  \filldraw[orange] (O1) circle (1.5pt) node[above right, text=orange] {$O_1$};
  \filldraw[orange] (O2) circle (1.5pt) node[left, text=orange] {$O_2$};

  % 旋转辅助线 (仅以顶点 A 的转变为例进行视觉指引)
  \draw[dashed, blue!50] (A) -- (A1);
  \draw[dashed, blue!50] (A1) -- (A');
  
  % O1 圆弧指示 (旋转 180度, 环绕 O1 的左侧)
  \draw[orange, ->] (O1) ++(100:0.8) arc (100:280:0.8);
  \node[orange, font=\small] at (-0.8, -1.7) {$180^\circ$};

  % O2 圆弧指示 (旋转 180度, 环绕 O2 的右下侧)
  \draw[orange, ->] (O2) ++(-112:0.8) arc (-112:68:0.8);
  \node[orange, font=\small] at (4.2, -1.4) {$180^\circ$};

\end{tikzpicture}
\end{center}

\vspace{0.4cm}
{\small\color{blue!70!black}
\textbf{绘图原理与步骤：}\\
\textbf{1. 证明拆解依据：} 数学性质定理指引：任意两次中心点的 $180^\circ$ 旋转必定等同于单一的一个平移变换，并且该平移向量 $\vec{v}$ 的模等同于两旋转中心距离 $O_1 O_2$ 的 2 倍（极向同理）。\\
\textbf{2. 虚拟构建系建立：} 设原 $\triangle ABC$ 左下角置于第三象限。平移出红色目标图形 $\triangle A'B'C'$。\\
\textbf{3. 推导出中转图与 O 点：} 任意选定点 $O_1(0.5, -1.5)$，并藉 $A \to A_1$ 推算全图点绕 $O_1$ 的 $180^\circ$（点反衍）坐标，汇聚成由灰色虚线构成的过渡态 $\triangle A_1B_1C_1$。随后严格定义出过渡点 $A_1$ 运动到终点 $A'$ 的中点即为确切的二次旋转圆心 $O_2(3,-1)$。\\
\textbf{4. 定向旋转诠释辅助线：} 不添加繁杂的全部网状连线，仅抽取最具有代表性的端点动作流 $A \to O_1 \to A_1 \to O_2 \to A'$，以蓝紫色虚点划线作为骨架关联，以贴身橙色 $180^\circ$ 弧形指向标展示连续两次扭转。
}

%% ==============================
%% 第66题（补充题）：选择合适的统计图
%% ==============================
\newpage



% ==============================
% 第67题（补充题）：三角形中心对称作图
% ==============================
\newpage

\section{解答：三角形中心对称的三种情境作图}

\vspace{0.6cm}

\textbf{\LARGE 58.} \textbf{（第 4 题）按下列条件，分别画一个与已知 $\triangle ABC$ 成中心对称的三角形：\\
（1）在图①中以顶点 $C$ 为对称中心；\\
（2）在图②中以 $AB$ 的中点 $M$ 为对称中心；\\
（3）在图③中以 $\triangle ABC$ 内的点 $P$ 为对称中心。}

\vspace{0.4cm}
\textbf{解：}

\textbf{画图如下：}

\begin{center}
  % 第一幅图：以顶点 C 为中心
  \begin{minipage}[b]{0.31\textwidth}
    \centering
    \begin{tikzpicture}[>=Stealth, scale=0.6]
      % 定义顶点
      \coordinate (B) at (0, 0);
      \coordinate (C) at (3, 0);
      \coordinate (A) at (1.2, 2.5);
      
      % 计算对称点 (中点C)
      \coordinate (A') at (4.8, -2.5); 
      \coordinate (B') at (6, 0);      
      \coordinate (C') at (3, 0);
      
      % 连线并标记对称路径
      \draw[dashed, blue!60] (A) -- (A');
      \draw[dashed, blue!60] (B) -- (B');
      
      % 原图形
      \draw[thick] (A) -- (B) -- (C) -- cycle;
      \filldraw (A) circle (1.5pt) node[above] {$A$};
      \filldraw (B) circle (1.5pt) node[above left] {$B$};
      \filldraw (C) circle (1.5pt) node[above right] {$C$};
      
      % 目标图形
      \draw[thick, red] (A') -- (B') -- (C') -- cycle;
      \filldraw[red] (A') circle (1.5pt) node[below, text=red] {$A'$};
      \filldraw[red] (B') circle (1.5pt) node[above right, text=red] {$B'$};
      % C' 和 C 重合，共同节点无需重复标注以免拥挤
      
      % 图名
      \node[below, font=\normalsize] at (3, -3.2) {①};
    \end{tikzpicture}
  \end{minipage}
  \hfill
  % 第二幅图：以中点 M 为中心
  \begin{minipage}[b]{0.31\textwidth}
    \centering
    \begin{tikzpicture}[>=Stealth, scale=0.6]
      \coordinate (B) at (0, 0);
      \coordinate (C) at (3, 0);
      \coordinate (A) at (1.2, 2.5);
      \coordinate (M) at (0.6, 1.25);
      
      % 计算对称点 (中点M)
      \coordinate (A') at (0, 0);
      \coordinate (B') at (1.2, 2.5);
      \coordinate (C') at (-1.8, 2.5);
      
      % 连线并标记对称路径
      \draw[dashed, blue!60] (C) -- (C');
      % A和B的路径重合在边AB上，为了图面干净不补充虚线
      
      % 原图形
      \draw[thick] (A) -- (B) -- (C) -- cycle;
      \filldraw (A) circle (1.5pt) node[above right] {$A$};
      \filldraw (B) circle (1.5pt) node[below right] {$B$};
      \filldraw (C) circle (1.5pt) node[right] {$C$};
      \filldraw (M) circle (1.5pt) node[above left] {$M$};
      
      % 目标图形
      \draw[thick, red] (A') -- (B') -- (C') -- cycle;
      \filldraw[red] (C') circle (1.5pt) node[left, text=red] {$C'$};
      \node[left, text=red] at (A') {$A'$};
      \node[above left, text=red] at (B') {$B'$};
      
      % 图名
      \node[below, font=\normalsize] at (0.6, -0.7) {②};
    \end{tikzpicture}
  \end{minipage}
  \hfill
  % 第三幅图：以内部点 P 为中心
  \begin{minipage}[b]{0.31\textwidth}
    \centering
    \begin{tikzpicture}[>=Stealth, scale=0.6]
      \coordinate (B) at (0, 0);
      \coordinate (C) at (3, 0);
      \coordinate (A) at (1.2, 2.5);
      \coordinate (P) at (1.5, 1.0);
      
      % 计算对称点 (中点P)
      \coordinate (A') at (1.8, -0.5);
      \coordinate (B') at (3, 2.0);
      \coordinate (C') at (0, 2.0);
      
      % 连线并标记对称路径
      \draw[dashed, blue!60] (A) -- (A');
      \draw[dashed, blue!60] (B) -- (B');
      \draw[dashed, blue!60] (C) -- (C');
      
      % 原图形
      \draw[thick] (A) -- (B) -- (C) -- cycle;
      \filldraw (A) circle (1.5pt) node[above] {$A$};
      \filldraw (B) circle (1.5pt) node[below left] {$B$};
      \filldraw (C) circle (1.5pt) node[below right] {$C$};
      \filldraw (P) circle (1.5pt) node[below right] {$P$};
      
      % 目标图形
      \draw[thick, red] (A') -- (B') -- (C') -- cycle;
      \filldraw[red] (A') circle (1.5pt) node[below right, text=red] {$A'$};
      \filldraw[red] (B') circle (1.5pt) node[right, text=red] {$B'$};
      \filldraw[red] (C') circle (1.5pt) node[left, text=red] {$C'$};
      
      % 图名
      \node[below, font=\normalsize] at (1.5, -1.2) {③};
    \end{tikzpicture}
  \end{minipage}
\end{center}

\vspace{0.4cm}
\textbf{作图说明：}

对于所有的中心对称作图，均遵循“\textbf{两点确定一线段，反向延长取等长}”的几何原则：
① 找出原图形的全部关键折点（即顶点 $A$、$B$、$C$）；
② 将每个顶点分别连接各自给定的对称中心（$C$点、$M$点或$P$点）；
③ 将该连线沿着中心点反向延长相等的长度，即可描出原顶点对应的中心对称点位置（即点 $A'$、$B'$、$C'$）；
④ 顺次使用直尺连接找出的所有对应对称点，即得到所求的中心对称图形 $\triangle A'B'C'$。

\vspace{0.4cm}
{\small\color{blue!70!black}
\textbf{绘图原理与步骤：}\\
\textbf{1. 结构统筹：} 考量到空间利用度，使用了 \texttt{minipage} 环境配合 \texttt{\textbackslash{}hfill} 均分出三个独立但等宽并列的 $0.31\backslash{}textwidth$ 画布区域；由于对称位置导致纵坐标大幅度偏移，巧妙利用了 \texttt{[b]} 底部对齐参数，分别单独抽离调校图编号①②③横向节点处于画面底端的特定相对 $y$ 坐标处（$-3.2$、$-0.7$、$-1.2$），从而令三张高低错落的解答图表在水平版面上完美实现物理“贴地基准线”对齐。\\
\textbf{2. 坐标同源定义：} 为了呈现直观的对比参照性，三幅子图均使用了一模一样原始姿态的 $\triangle ABC$：选定参数为 $B(0,0)$、$C(3,0)$、$A(1.2, 2.5)$。\\
\textbf{3. 中点射线拓展：} 遵循“点关于原点中心对称”等价于“两点中点公式倍积”的硬逻辑 $X'=2 \cdot O-X$。依次按子题要求将标靶参考点 $O$ 置为 $C(3,0)$、$M(0.6, 1.25)$、$P(1.5, 1.0)$，并利用脚本运算得到所有红色的全映射点。顺带一提，在图②中，按数学定义推算出的 $A'$ 与 $B'$ 十分类似于在边上原地互换。\\
\textbf{4. 反差显隐：} 加粗的黑色墨水线构筑母体，对比强烈的红色彩线构筑对碰的本体，最后辅助串通原与象点的长跨度蓝紫色虚参考线，以最高程度保留了辅助理解思考的画法痕迹。
}

%% ==============================
%% 第68题（补充题）：网格中的平移与旋转/中心对称综合
%% ==============================
\newpage



%% ==============================
%% 第68题（补充题）：网格中的平移与旋转/中心对称综合
%% ==============================
\newpage

\section{解答：网格图中的平移、旋转与中心对称}

\vspace{0.6cm}

\textbf{\LARGE 59.} \textbf{（第 7 题）如图，在方格纸中，$\triangle A_1B_1C_1$ 和 $\triangle A_2B_2C_2$ 是两个形状、大小都一样的三角形。\\
（1）$\triangle A_1B_1C_1$ 经过怎样的平移、旋转变换后，可以与 $\triangle A_2B_2C_2$ 重合？\\
（2）$\triangle A_1B_1C_1$ 经过怎样的变换后，可以与 $\triangle A_2B_2C_2$ 成中心对称？画出变换后的三角形，并标出对称中心。}

\vspace{0.4cm}
\textbf{解：}

\textbf{（1）}答案不唯一，参考变幻路径如下：\\
将 $\triangle A_1B_1C_1$ 先\textbf{向上}平移 $4$ 个单位长度，再\textbf{向右}平移 $3$ 个单位长度；再绕平移后点 $C_1$ 的对应点 $C_1'$ 按\textbf{顺时针}方向旋转 $90^\circ$ 即可重合于 $\triangle A_2B_2C_2$。（\textit{注：原参考答案笔误为“逆时针”，从网格坐标可严格计算向量推演应为顺时针旋转。}）

\vspace{0.3cm}
\textbf{（2）}答案不唯一，参考变换路径如下：\\
将 $\triangle A_1B_1C_1$ 绕点 $C_1$ 按\textbf{逆时针}方向旋转 $90^\circ$ 即可得到变换后的目标三角形 $\triangle A_3B_3C_3$。此时的 $\triangle A_3B_3C_3$ 与 $\triangle A_2B_2C_2$ 成中心对称关系。对称中心 $O$ 即为线段 $C_3C_2$ 的中点。

\textbf{画图如下：}

\begin{center}
\begin{tikzpicture}[>=Stealth, scale=0.8]
  % 绘制网格背景 (0 到 13, 0 到 9)
  \draw[dashed, gray] (0,0) grid (13,9);
  
  % 定义坐标
  \coordinate (A1) at (4, 5);
  \coordinate (B1) at (3, 2);
  \coordinate (C1) at (5, 2);
  
  % 修正 A2 为 11 以确保形状大小完全一致
  \coordinate (A2) at (11, 7);
  \coordinate (B2) at (8, 8);
  \coordinate (C2) at (8, 6);
  
  % 第(1)问变换中间状态 A1'B1'C1' (上移4，右移3)
  \coordinate (A1p) at (7, 9);
  \coordinate (B1p) at (6, 6);
  \coordinate (C1p) at (8, 6);

  % 第(2)问变换后的三角形 A3B3C3 (A1B1C1 绕 C1 逆时针单中心旋转 90度)
  \coordinate (A3) at (2, 1);
  \coordinate (B3) at (5, 0);
  \coordinate (C3) at (5, 2); % 与 C1 重合
  
  % 对称中心 O，为 C3(5,2) 和 C2(8,6) 的中点 -> (6.5, 4)
  \coordinate (O) at (6.5, 4);

  % 1. 画中间状态 A1'B1'C1' (灰色虚线)
  \draw[dashed, gray, thick] (A1p) -- (B1p) -- (C1p) -- cycle;
  \node[above, text=gray] at (A1p) {$A_1'$};
  \node[left, text=gray] at (B1p) {$B_1'$};
  \node[above left, text=gray] at (C1p) {$C_1'$};
  
  % 画平移箭头指示
  \draw[->, dotted, thick, gray] (A1) -- ++(0,4) node[midway, left, font=\scriptsize] {上移 4} -- (A1p) node[midway, above, font=\scriptsize] {右移 3};
  
  % 画第一问的顺时针旋转指示
  % C1p(8,6), A1p(7,9) -> angle = 108.4度
  % C2(8,6), A2(11,7) -> angle = 18.4度
  \draw[->, orange, thick] (C1p) ++(108.4:1.2) arc (108.4:18.4:1.2);
  \node[orange, font=\small] at (8.7, 7.5) {$90^\circ$顺转};

  % 2. 绘制原三角形 A1B1C1
  \draw[thick] (A1) -- (B1) -- (C1) -- cycle;
  \filldraw[black] (A1) circle (1.5pt) node[above right] {$A_1$};
  \filldraw[black] (B1) circle (1.5pt) node[below left] {$B_1$};
  % C1 统一在下面标记

  % 3. 绘制目标三角形 A2B2C2
  \draw[thick] (A2) -- (B2) -- (C2) -- cycle;
  \filldraw[black] (A2) circle (1.5pt) node[right] {$A_2$};
  \filldraw[black] (B2) circle (1.5pt) node[above left] {$B_2$};
  % C2 统一在下面标记

  % 4. 绘制第(2)问的三角形 A3B3C3 (红色实线)
  \draw[thick, red] (A3) -- (B3) -- (C3) -- cycle;
  \filldraw[red] (A3) circle (1.5pt) node[left, text=red] {$A_3$};
  \filldraw[red] (B3) circle (1.5pt) node[below, text=red] {$B_3$};
  
  % C1/C3 重合节点统一标记
  \filldraw[black] (C1) circle (1.5pt);
  \node[below right] at (C1) {$C_1(C_3)$};
  
  % C1p/C2 重合节点统一标记
  \filldraw[black] (C2) circle (1.5pt);
  \node[below right] at (C2) {$C_2$};

  % 5. 绘制中心对称辅助线与中心 O
  \draw[dashed, blue!50] (A3) -- (A2);
  \draw[dashed, blue!50] (B3) -- (B2);
  \draw[dashed, blue!50] (C3) -- (C2);
  
  \filldraw[blue] (O) circle (1.5pt) node[above=2pt, text=blue] {$O$};

\end{tikzpicture}
\end{center}

\vspace{0.4cm}
{\small\color{blue!70!black}
\textbf{绘图原理与步骤：}\\
\textbf{1. 纠错溯源法（第一问）：} 建立 $13 \times 9$ 规格方格纸系统。依据图像比例“两个形状大小都一样”，将输入的坐标修正为 $A_1(4,5)$、$B_1(3,2)$、$C_1(5,2)$ 以及 $A_2(\mathbf{11},7)$、$B_2(8,8)$、$C_2(8,6)$，纠正了原题干文本错漏的 $x=14$ 以保障绝对全等。严谨推算两原图形态：平移至过渡点并形成灰色 $\triangle A_1'B_1'C_1'$ 后，由于内部矢量方向从水平向右的 $(2,0)$ 被映射为了竖直向下的 $(0,-2)$，此显然为典型的“\textbf{顺时针 $\mathbf{90^\circ}$ 旋转}”，通过橙色弧线指示并在题解中直接纠正了原参考答案的笔误。\\
\textbf{2. 方位倒推（第二问）：} 因为 $\triangle A_2B_2C_2$ 是原图顺时针转了 $90^\circ$；要得到能与之发生中心对称（即倒置 $180^\circ$ 差异）的新三角形，反向推导只需将原图绕 $C_1$ \textbf{逆时针转 $\mathbf{90^\circ}$} 即可恰好满足角度咬合关联。以此确算出全图三点的各自新落点 $A_3(2,1)$、$B_3(5,0)$、$C_3(5,2)$，并标红显示。\\
\textbf{3. 中点定核：} 取 $C_3(5,2)$ 与 $C_2(8,6)$ 的线段计算出严格数学中点 $O(6.5,4)$。随后施加跨图形的极坐标对角蓝紫虚线，完美匹配每一对相关联红色端点与黑色端点均穿过该共点 $O$，以强有力的视觉效果成功验证并突写了中心对称性质规律！
}

%% ==============================
%% 第69题（补充题）：全等三角形的图形变换
%% ==============================
\newpage



%% ==============================
%% 第69题（补充题）：全等三角形的图形变换
%% ==============================
\newpage

\section{解答：全等三角形的图形变换}

\vspace{0.6cm}

\textbf{\LARGE 60.} \textbf{（例 2）如图 9-23，方格纸中有两个形状、大小都相同的三角形，通过怎样的图形变换可以使其中一个三角形与另一个三角形重合？请描述图形变换的过程。}

\vspace{0.4cm}
\textbf{解：}

两个三角形均为直角三角形，直角边分别为 $2$ 和 $3$（斜边 $= \sqrt{13}$），形状大小完全相同。

\textbf{图形如下：}

\begin{center}
\begin{tikzpicture}[>=Stealth, scale=1]
  % 绘制网格背景 (0 到 8, 0 到 6)
  \draw[dashed, gray] (0,0) grid (8,6);
  
  %% ===== 左侧三角形（直角在 B） =====
  \coordinate (A) at (1, 2);
  \coordinate (B) at (3, 2);
  \coordinate (C) at (3, 5);
  
  % 绘制三角形
  \draw[thick] (A) -- (B) -- (C) -- cycle;
  \filldraw (A) circle (1.2pt) node[below left] {$A$};
  \filldraw (B) circle (1.2pt) node[below right] {$B$};
  \filldraw (C) circle (1.2pt) node[above] {$C$};
  
  % % 直角标记 (B 处)
  % \draw[thin] (2.75, 2) -- (2.75, 2.25) -- (3, 2.25);
  
  %% ===== 右侧三角形（直角在 E, 非格点坐标 76/13, 49/13） =====
  \coordinate (D) at (4, 3);
  \coordinate (E) at ({76/13}, {49/13});
  \coordinate (F) at (7, 1);
  
  % 绘制三角形
  \draw[thick] (D) -- (E) -- (F) -- cycle;
  \filldraw (D) circle (1.2pt) node[above left] {$D$};
  \filldraw (E) circle (1.2pt) node[above right] {$E$};
  \filldraw (F) circle (1.2pt) node[below right] {$F$};
  
  % % 直角标记 (E 处)
  % % ED 单位向量 (-12/13, -5/13), EF 单位向量 (5/13, -12/13)
  % \coordinate (Eq1) at ($(E) + ({-12/13*0.25}, {-5/13*0.25})$);
  % \coordinate (Eq2) at ($(E) + ({(-12+5)/13*0.25}, {(-5-12)/13*0.25})$);
  % \coordinate (Eq3) at ($(E) + ({5/13*0.25}, {-12/13*0.25})$);
  % \draw[thin] (Eq1) -- (Eq2) -- (Eq3);

  %% ===== 变换过程可视化 =====
  
  % 步骤①：平移后的中间三角形 A'B'C'（灰色虚线）
  \coordinate (Ap) at (4, 3);   % A→A'=D
  \coordinate (Bp) at (6, 3);
  \coordinate (Cp) at (6, 6);
  \draw[dashed, blue!70, thick] (Ap) -- (Bp) -- (Cp) -- cycle;
  \node[below, text=blue!60, font=\small] at (Bp) {$B'$};
  \node[above, text=blue!60, font=\small] at (Cp) {$C'$};
  
  % 平移箭头 A→D (dotted)
  \draw[->, dotted, thick, blue!60] (A) -- (D);
  \node[above left, text=blue!60, font=\scriptsize] at (2.6, 1.9) {平移};

  % 步骤②：反射对称轴（过 D，斜率 1/5 的绿色虚线）
  % 直线方程 y - 3 = (1/5)(x - 4)，即 x=0→y=2.2, x=8→y=3.8
  \draw[densely dashdotted, green!60!black, thick] (0, {3-4/5}) -- (8, {3+4/5});
  \node[above, text=green!60!black, font=\scriptsize] at (0.8, {3-4/5+0.1}) {对称轴};

  % 反射辅助线 B'→E 和 C'→F（蓝色虚线）
  \draw[dashed, green!60!black] (Bp) -- (E);
  \draw[dashed, green!60!black] (Cp) -- (F);

\end{tikzpicture}
\end{center}

\vspace{0.3cm}
\textbf{变换过程描述：}

答案不唯一，以下为一种参考变换过程：

\textbf{①} 先将左侧 $\triangle ABC$ \textbf{向右平移 $\mathbf{3}$ 个单位}、\textbf{向上平移 $\mathbf{1}$ 个单位}，使顶点 $A(1,2)$ 与右侧三角形的顶点 $D(4,3)$ 重合；

\textbf{②} 再将平移后的三角形关于过点 $D$ 的一条对称轴（该轴斜率为 $\dfrac{1}{5}$）作\textbf{轴对称变换}，即可与 $\triangle DEF$ 完全重合。

\vspace{0.2cm}
{\small
\textbf{验证：} 平移后三角形顶点为 $A'(4,3)$, $B'(6,3)$, $C'(6,6)$。将 $B'(6,3)$ 关于直线 $x - 5y + 11 = 0$ 作对称得 $(76/13, 49/13) = E$ $\checkmark$；将 $C'(6,6)$ 作对称得 $(7, 1) = F$ $\checkmark$。
}

\vspace{0.4cm}
{\small\color{blue!70!black}
\textbf{绘图原理与步骤：}\\
\textbf{1. 网格系统：} 建立 $8 \times 6$ 规格方格纸，$x \in [0,8]$，$y \in [0,6]$。\\
\textbf{2. 坐标设定：} 左侧三角形 $A(1,2)$、$B(3,2)$、$C(3,5)$，直角在 $B$。右侧三角形 $D(4,3)$、$F(7,1)$ 为格点；第三顶点 $E$ 为非格点，由 $|DE|=2$、$|EF|=3$、$DE \perp EF$ 联立方程组解得精确坐标 $E = (76/13,\; 49/13) \approx (5.85,\, 3.77)$。\\
\textbf{3. 全等验证：} 两三角形直角边均为 $2$ 和 $3$，斜边均为 $\sqrt{13}$，满足 SSS 全等。\\
\textbf{4. 变换分解：} 平移向量 $(3,1)$ 使 $A \to D$。平移后 $B' \to (6,3)$、$C' \to (6,6)$，需映射至 $E$、$F$。变换矩阵 $M = \bigl[\frac{12}{13},\, \frac{5}{13};\; \frac{5}{13},\, {-}\frac{12}{13}\bigr]$，$\det M = -1$，为关于过 $D$ 斜率 $\frac{1}{5}$ 直线的轴对称。\\
\textbf{5. 直角标记：} $B$ 处沿坐标轴方向绘制。$E$ 处沿 $ED$、$EF$ 单位方向各取 $0.25$ 步长构成倾斜直角方块。
}

%% ==============================
%% 第70题（补充题）：不等式的数轴表示
%% ==============================


%% ==============================
%% 第87题（补充题）：三角形旋转与平行四边形判定
%% ==============================
\newpage

\section{解答：图形旋转与四边形判定}

\vspace{0.6cm}

\textbf{\LARGE 78.} \textbf{如图 $8-5$，把 $\triangle ABC$ 绕点 $C$ 旋转 $180^\circ$ 得 $\triangle DEC$，请画出图形并判断四边形 $ABDE$ 的形状。}

\vspace{0.4cm}
\textbf{解：}

\vspace{0.2cm}
\begin{center}
\begin{tikzpicture}[>=Stealth, x=1cm, y=1cm]
  % 定义点 C 作为旋转中心（原点）
  \coordinate (C) at (0, 0);
  % 根据原图比例估算 A 和 B 的相对位置
  \coordinate (B) at (-3, 0);
  \coordinate (A) at (-0.8, 2);
  
  % 根据绕点 C 旋转 180° 的性质求 D 和 E 的坐标
  % D 是 A 的对应点，E 是 B 的对应点
  \coordinate (D) at (0.8, -2);
  \coordinate (E) at (3, 0);
  
  % 1. 原图形：绘制 \Delta ABC（黑色实线）
  \draw[line width=0.8pt] (A) -- (B) -- (C) -- cycle;
  \filldraw[black] (A) circle (1.5pt) node[above, font=\small] {$A$};
  \filldraw[black] (B) circle (1.5pt) node[below left, font=\small] {$B$};
  \filldraw[black] (C) circle (1.5pt) node[above right, font=\small] {$C$};
  
  % 2. 旋转后的图形：绘制 \Delta DEC（蓝色实线）
  \draw[line width=0.8pt, blue!80!black] (D) -- (E) -- (C) -- cycle;
  \filldraw[black] (D) circle (1.5pt) node[below, font=\small, blue!80!black] {$D$};
  \filldraw[black] (E) circle (1.5pt) node[above right, font=\small, blue!80!black] {$E$};
  
  % 3. 连接构成的四边形 ABDE
  % 边 AB 和 DE 已经在三角形中画出，这里补充画边 BD 和 AE，使用虚线
  \draw[line width=0.8pt, dashed, red!80!black] (A) -- (E);
  \draw[line width=0.8pt, dashed, red!80!black] (B) -- (D);
  
\end{tikzpicture}
\end{center}

\vspace{0.4cm}
\textbf{判断结果：四边形 $\boldsymbol{ABDE}$ 是平行四边形。}

\textbf{理由：}
根据旋转 $180^\circ$（即中心对称）的性质可知，点 $A, C, D$ 在同一直线上且 $AC = CD$，说明点 $C$ 是线段 $AD$ 的中点；同理，点 $B, C, E$ 在同一直线上且 $BC = CE$，说明点 $C$ 也是线段 $BE$ 的中点。

根据“\textbf{对角线互相平分的四边形是平行四边形}”这一判定定理，因为四边形 $ABDE$ 的对角线 $AD$ 和 $BE$ 在点 $C$ 处互相平分，所以四边形 $ABDE$ 是平行四边形。

\vspace{0.4cm}
{\small\color{blue!70!black}
\textbf{绘图原理与步骤：}\\
\textbf{1. 锚点与旋转代数求解：} 将原像点 $C$ 设为直角坐标系的原点 $(0,0)$。根据图像估算置点 $B$ 于水平轴上 $(-3,0)$，置点 $A$ 于 $(-0.8, 2)$ 处。根据中心对称（即绕点 $C$ 旋转 $180^\circ$）规则，直接取横纵坐标相反数，计算得对应点 $D(0.8, -2)$ 与 $E(3, 0)$，实现代数级别的精准刚体旋转。\\
\textbf{2. 层次化元素图解：} 采用粗黑实线重绘原始结构 $\triangle ABC$ 作为主从视觉基准。随后以醒目的蓝色实线勾勒出经过中心对称重构的 $\triangle DEC$ 的像。\\
\textbf{3. 四边形闭合及证明铺垫：} 利用红色的虚线轻巧补齐待判断四边形的外围缺口轮廓（边 $AE$ 和 $BD$）。这样的色彩与虚实搭配不仅完整呈现了最终生成的四边形闭包实体，同时在视觉上保留并凸显了主轴线 $AD$ 与 $BE$ “既是两个独立三角形的边、又是统摄四边形的交叉对角线”的身份，从而在直觉层面不证自明地诠释了对角线平分的判定逻辑。
}

%% ==============================
%% 第88题（补充题）：在方格图中拼图
%% ==============================
\newpage


